\section{Supporting Experimental Results}
\label{app:retained-results}
Paired intervals follow Appendix~\ref{app:independent-readouts}.

\subsection{Controlled correction and target choice}
\label{app:closed-form-controlled}
\label{app:target-controls}
\paragraph{Closed-form correction.}
In Table~\ref{tab:controlled-cf}, Closed-form MGPA changes source AUROC from
identity by $-.173$ $[-.180,-.164]$. Frozen- and refitted-task changes are
$-.0002$ $[-.0007,+.0002]$ and $+.0004$ $[-.0001,+.0008]$;
the refitted-task gain over LEACE is $+.0879$ $[.0762,.0998]$.

\paragraph{Conditional versus unconditional targets.}
The unconditional control uses each critic's empirical mean score on the
current corrected FIT inputs, recomputed at every stage without downstream
labels or HEAD/EVAL data. Gate, critic fitting, solver and the 48-stage
horizon are unchanged. Each trajectory refits its critics, testing learned
correction rather than the fixed-score equivalence of
Proposition~\ref{prop:parallel-levels}. Zero-score targeting remains a
separate control: balanced source labels do not imply a zero mean logit.

\begin{table}[htbp]
\centering\small
\caption{\textbf{Conditional versus constant targets.}
Native controlled outputs for the same 81 participants and fits as
Figure~\ref{fig:target-cost}c. $S$: source AUROC; $T_f/T_r$: frozen/refitted
task AUROC; $M$: movement.}
\label{tab:target-controls}
% Generated by build_experiment_assets.py; see experiment_provenance.json.
\begingroup\small\setlength{\tabcolsep}{4pt}
\begin{tabular}{@{}lrrrr@{}}
\toprule
Target & $S\downarrow$ & $T_f\uparrow$ & $T_r\uparrow$ & $M$ \\
\midrule
Conditional mean & 0.5105 & 0.7496 & 0.7527 & 0.1500 \\
Unconditional mean (FIT) & 0.5214 & 0.7488 & 0.7574 & 0.4063 \\
Zero score & 0.5199 & 0.7479 & 0.7562 & 0.4172 \\
\bottomrule
\end{tabular}\endgroup

\end{table}

Conditional minus empirical-unconditional anchoring gives movement
$-.2563$ $[-.2907,-.2232]$ and source AUROC $-.0110$ $[-.0143,-.0079]$.
Frozen-task change is $+.0008$ $[-.0015,.0030]$; refitted-task change is
$-.0047$ $[-.0070,-.0025]$.

\subsection{Temporal N170: nonlinear correction and task utility}
\label{app:critic-complexity}
\label{app:igbp-comparison}
\paragraph{Linear-only ablation.}
Removing the GELU critic retains the gate, Ridge mean anchor, smooth
weighting, trust-region radius and stage horizon. Single-critic damping
$.0005$ preserves regularization relative to mean critic loss, equivalently
duplicating the affine constraint at the shared two-critic damping.
Sample-dependent bounded steps make this a linear-critic ablation, not a
globally affine map. Closed-form MGPA instead estimates its direction and
covariance metric from measurement contrasts and the joint residual model.
Table~\ref{tab:critic-complexity} shows why the distinction matters:
linear-only correction lowers linear accessibility but leaves nonlinear
accessibility nearly intact. Iterative MGPA lowers the full-bank score
with higher task scores at nearly the same movement.

\begin{table}[htbp]
\centering\small
\caption{\textbf{Nonlinear critics on N170.} A: strongest linear reader $L$
versus full-bank $S$ for the ten participants in Table~\ref{tab:real-main}A.
B: Iterative minus linear-only AUROC percentage points, with paired 95\%
intervals. Iterative estimates average fits; linear-only is deterministic.}
\label{tab:critic-complexity}
\label{tab:critic-complexity-paired}
% Generated by build_experiment_assets.py; see experiment_provenance.json.
\begingroup\small\setlength{\tabcolsep}{4pt}
\begin{tabular*}{\textwidth}{@{\extracolsep{\fill}}lrr@{}}
\toprule
\multicolumn{3}{@{}l}{\textbf{A. Native source accessibility}} \\
Correction & $L\downarrow$ & $S\downarrow$ \\
\midrule
Identity & 0.763 & 0.989 \\
MGPA-Iter & 0.794 & 0.946 \\
Linear-only & 0.663 & 0.987 \\
\bottomrule
\end{tabular*}\endgroup
\par\vspace{4pt}
% Generated by build_experiment_assets.py; see experiment_provenance.json.
\begingroup\small\setlength{\tabcolsep}{4pt}
\begin{tabular*}{\textwidth}{@{\extracolsep{\fill}}lrrr@{}}
\toprule
\multicolumn{4}{@{}l}{\textbf{B. MGPA-Iter minus Linear-only (AUROC percentage points)}} \\
Contrast & $\Delta S\downarrow$ [95\% CI] & $\Delta T_f\uparrow$ [95\% CI] & $\Delta T_r\uparrow$ [95\% CI] \\
\midrule
Paired difference & -4.13 [-6.47, -2.10] & +0.93 [-0.05, +1.88] & +1.71 [+0.05, +3.84] \\
\bottomrule
\end{tabular*}\endgroup

\end{table}

\paragraph{Comparison with IGBP.}
At native outputs, Iterative MGPA minus IGBP gives frozen/refitted task
gains of $+.0123$ $[.0042,.0203]$ and $+.0361$ $[.0306,.0405]$, and movement
lower by $.1504$ $[.1230,.1769]$. The source-AUROC difference is
$-.0091$ $[-.0198,.0006]$. At common VAL movement budget $.25$
(Table~\ref{tab:igbp-budget}), the respective source, frozen-task and
refitted-task differences are $-.0234$ $[-.0393,-.0099]$,
$+.0082$ $[.0015,.0154]$ and $+.0156$ $[.0063,.0269]$.
At budget $.10$, source accessibility is lower but a downstream advantage
is not established. These comparisons assess complete correction
procedures; equal iteration counts would not equalize their computation.

\begin{table}[htbp]
\centering\small
\caption{\textbf{N170: MGPA versus IGBP at matched VAL budgets.}
First-crossing endpoints follow Appendix~\ref{app:gate-attribution};
attained EVAL movement $M$ may differ. Same ten participants and shared fits.}
\label{tab:igbp-budget}
% Generated by build_experiment_assets.py; see experiment_provenance.json.
\begingroup\small\setlength{\tabcolsep}{4pt}
\begin{tabular}{@{}llrrrr@{}}
\toprule
VAL budget & Method & $S\downarrow$ & $T_f\uparrow$ & $T_r\uparrow$ & $M$ \\
\midrule
.10 & MGPA-Iter & 0.9850 & 0.7274 & 0.7300 & 0.0974 \\
.10 & IGBP & 0.9893 & 0.7292 & 0.7284 & 0.0898 \\
\addlinespace[2pt]
.25 & MGPA-Iter & 0.9643 & 0.7292 & 0.7394 & 0.2422 \\
.25 & IGBP & 0.9877 & 0.7210 & 0.7237 & 0.2239 \\
\bottomrule
\end{tabular}\endgroup

\end{table}

\paragraph{Task utility.}
Iterative MGPA improves frozen/refitted task AUROC over identity by
$+.0070$ $[+.0012,+.0131]$ and $+.0187$ $[+.0042,+.0366]$, respectively.
Both transfer directions receive equal weight:
the refitted gain is concentrated in Neuroscan-to-Flex ($.608$ to $.648$),
while the reverse direction changes from $.847$ to $.844$.

\subsection{Recorded SSVEP: full-input accessibility and utility}
\label{app:recorded-full-input}
On the same saved outputs and participant roles, Iterative MGPA gives
pooled source AUROC $.7954$ and full-bank $S=.8718$, versus identity's
$.9155$. Readers trained directly on all tokens independently support
attenuation: their maximum falls from $.8987$ to $.8685$, a paired change
of $-.0302$ $[-.0459,-.0174]$. The full-bank difference from LEACE is
$-.0463$ $[-.0671,-.0225]$. Closed-form MGPA gives $S=.8949$, a change
of $-.0206$ $[-.0356,-.0043]$ from identity. All maps share the 41-reader bank
(Appendix~\ref{app:full-token-readers}).

Residual-only readers attain $.8718$ for every common-offset correction;
this determines the full-bank maximum for Iterative MGPA, linear-only MGPA
and IGBP. Thus the maximum also reflects retained temporal structure.
CORAL reaches $S=1$ in both orientations under the common-offset interface;
these are not results for token-specific CORAL maps.

\paragraph{Comparison with IGBP and task utility.}
IGBP and Iterative MGPA both attain full-bank $S=.8718$, with pooled-mean
scores $.7832$ and $.7954$, respectively. IGBP has lower movement
($.0956$ versus $.1490$), slightly higher frozen-task AUROC ($.6891$ versus
$.6880$), and similar refitted-task AUROC ($.6884$ versus $.6886$).
Against identity, Iterative MGPA's frozen/refitted task changes are
$-.0002$ $[-.0017,+.0013]$ and $+.0004$ $[-.0028,+.0034]$.

\subsection{P300-selected reuse: absolute outcomes and paired comparisons}
\label{app:transfer-results}
\label{app:reuse-results}
Tables~\ref{tab:reuse-absolute}--\ref{tab:reuse-paired} supplement the main
reuse comparison at its unchanged P300-selected operating points.
Both MGPA corrections improve
transferred worst AUROC over LEACE on N170 and MMN, with paired intervals
excluding zero; on P300 these intervals contain zero. Control-head changes
from identity are small: Closed-form MGPA gains $.0012$ $[.0006,.0018]$ on
MMN, while the other three MGPA transfer-task intervals contain zero.
N170 worst AUROC remains below .5: the gain is improved robustness rather
than complete recovery.

\begin{table}[htbp]
\centering\small
\caption{\textbf{Absolute reuse outcomes.} A/I/R: aligned/independent/reversed
AUROC of the shortcut-sensitive head; worst/control results are in
Table~\ref{tab:reuse-main}. Iterative entries average fits; other maps are
evaluated once. CORAL/FEATMAP require device identity.}
\label{tab:reuse-absolute}
% Generated by build_experiment_assets.py; see reuse_absolute_compact.provenance.json.
\begingroup\small\setlength{\tabcolsep}{4pt}
\begin{tabular*}{\textwidth}{@{\extracolsep{\fill}}lrrrrrrrrr@{}}
\toprule
 & \multicolumn{3}{c}{P300} & \multicolumn{3}{c}{N170} & \multicolumn{3}{c}{MMN} \\
\cmidrule(lr){2-4}\cmidrule(lr){5-7}\cmidrule(lr){8-10}
Method & A & I & R & A & I & R & A & I & R \\
\midrule
Identity & 0.770 & 0.568 & 0.355 & 0.914 & 0.547 & 0.122 & 0.753 & 0.603 & 0.446 \\
LEACE & 0.604 & 0.602 & 0.601 & 0.782 & 0.590 & 0.384 & 0.614 & 0.633 & 0.652 \\
MGPA-CF & 0.643 & 0.621 & 0.599 & 0.735 & 0.579 & 0.420 & 0.642 & 0.644 & 0.646 \\
MGPA-Iter & 0.636 & 0.617 & 0.599 & 0.744 & 0.594 & 0.441 & 0.632 & 0.639 & 0.646 \\
CORAL & 0.662 & 0.607 & 0.550 & 0.818 & 0.564 & 0.277 & 0.639 & 0.632 & 0.624 \\
FEATMAP & 0.625 & 0.600 & 0.573 & 0.697 & 0.591 & 0.471 & 0.623 & 0.621 & 0.618 \\
\bottomrule
\end{tabular*}\endgroup

\end{table}

\begin{table}[htbp]
\centering\small
\caption{\textbf{Paired reuse gains.} Method-minus-comparator AUROC and 95\%
intervals, recomputing worst AUROC inside each bootstrap draw
(Appendix~\ref{app:reuse-protocol}).}
\label{tab:reuse-paired}
% Generated by build_experiment_assets.py; see experiment_provenance.json.
\begingroup\small\setlength{\tabcolsep}{4pt}
\begin{tabular*}{\textwidth}{@{\extracolsep{\fill}}lllrr@{}}
\toprule
Task & Method & Reference & $\Delta$ worst [95\% CI] & $\Delta$ control head [95\% CI] \\
\midrule
P300 & MGPA-CF & Identity & $+0.244\ [+0.220, +0.268]$ & $+0.0005\ [-0.0010, +0.0022]$ \\
P300 & MGPA-CF & LEACE & $-0.002\ [-0.022, +0.027]$ & $+0.0050\ [+0.0036, +0.0066]$ \\
P300 & MGPA-Iter & Identity & $+0.244\ [+0.224, +0.262]$ & $+0.0003\ [-0.0016, +0.0025]$ \\
P300 & MGPA-Iter & LEACE & $-0.002\ [-0.016, +0.021]$ & $+0.0047\ [+0.0034, +0.0062]$ \\
N170 & MGPA-CF & Identity & $+0.298\ [+0.244, +0.340]$ & $-0.0014\ [-0.0044, +0.0019]$ \\
N170 & MGPA-CF & LEACE & $+0.036\ [+0.005, +0.067]$ & $+0.0002\ [-0.0023, +0.0037]$ \\
N170 & MGPA-Iter & Identity & $+0.319\ [+0.251, +0.380]$ & $-0.0010\ [-0.0035, +0.0013]$ \\
N170 & MGPA-Iter & LEACE & $+0.057\ [+0.034, +0.083]$ & $+0.0006\ [-0.0009, +0.0025]$ \\
MMN & MGPA-CF & Identity & $+0.195\ [+0.168, +0.215]$ & $+0.0012\ [+0.0006, +0.0018]$ \\
MMN & MGPA-CF & LEACE & $+0.028\ [+0.002, +0.034]$ & $+0.0024\ [+0.0006, +0.0043]$ \\
MMN & MGPA-Iter & Identity & $+0.186\ [+0.157, +0.209]$ & $+0.0003\ [-0.0014, +0.0017]$ \\
MMN & MGPA-Iter & LEACE & $+0.019\ [+0.004, +0.024]$ & $+0.0016\ [-0.0005, +0.0036]$ \\
\bottomrule
\end{tabular*}\endgroup

\end{table}

Across the three Iterative MGPA fits, worst-association AUROC ranges from
$.417$ to $.471$ on N170 and $.622$ to $.639$ on MMN. Every P300-selected map,
reused unchanged, improves this score over LEACE on both transfer tasks;
Table~\ref{tab:reuse-paired} reports intervals for the mean across fits.
\FloatBarrier
